\section{\uppercase{Introduction}}
\label{sec:introduction}

Multi-temporal Building Change Detection (BCD) from very high–resolution (VHR) satellite imagery plays a crucial role in urban expansion monitoring, infrastructure management, and post–disaster assessment \cite{Daudt2018SiamICIP,Zhang2020FusionNet,Hussain2013Review}. Given two or more images of the same geographic area acquired at different times, the goal is to determine whether individual buildings have been newly constructed, removed, or altered in their footprint or roof structure. The task is typically formulated as a bi-temporal semantic segmentation problem \cite{Zhang2020Survey,Shafique2022Review}, where a model predicts a binary change mask from paired inputs. Public benchmarks such as LEVIR-CD \cite{Zheng2020LEVIR} and DSIFN-CD \cite{Shen2021DSIFN} have been instrumental in advancing BCD because they offer sufficiently large, diverse, and well-annotated bi-temporal samples to train and reliably evaluate deep models on building-scale structural changes. Fundamentally, accurately predicting change regions requires a dual capability, the model must: (i) resolve subtle, fine–grained alterations at the building level, and (ii) integrate broader spatial–temporal context to disambiguate true structural changes from irrelevant variations such as illumination or seasonal effects. In this context, RALA-ChangeFormer \cite{Fan2025RALA} satisfies the dual requirement while substantially reducing computational cost in comparison to the baseline ChangeFormer \cite{Bandara2022ChangeFormer}. Table~\ref{tab:intro_teaser} shows a comparison of RALA-ChangeFormer and ChangeFormer with different datasets.

\begin{table}[t]
\centering
\caption{Comparison between ChangeFormer (CF) and the proposed RALA-ChangeFormer (RALA-CF) on two benchmarks with input dimension 256 x 256 and batch 16.}
\label{tab:intro_teaser}
\resizebox{\columnwidth}{!}{
        \begin{tabular}{lcc}
        \hline
        \textbf{Dataset} & \textbf{ChangeFormer} & \textbf{RALA-CF} \\
        \hline
        \multicolumn{3}{l}{\textbf{LEVIR-CD \cite{Zheng2020LEVIR}}} \\[2pt]
        F1          & 90.40  & 89.87 \\
        IoU         & 82.48  & 81.60 \\
        FLOPs (G)   & 138.77 & \textbf{135.62} \\
        \hline
        \multicolumn{3}{l}{\textbf{DSIFN-CD \cite{Shen2021DSIFN}}} \\[2pt]
        F1          & 86.67  & 84.07 \\
        IoU         & 76.48  & 73.96 \\
        FLOPs (G)   & 248.77 & \textbf{228.62} \\
        \hline
        \end{tabular}
}
\end{table}

In the same vein, motivated by the challenges inherent to bi-temporal semantic segmentation, recent years have witnessed a shift from convolutional Siamese architectures to Transformer-based models for BCD. Methods such as BIT \cite{Chen2022BIT}, SNUNet \cite{Fang2021SNUNetCD} and ChangeFormer \cite{Bandara2022ChangeFormer} leverage self-attention to capture long-range dependencies and global spatial–temporal context, achieving state-of-the-art performance on multiple datasets. Vision Transformer (ViT) architectures \cite{Dosovitskiy2021ViT,Raghu2023RevengeViT, Fan2024ViTAR} have demonstrated strong representational capacity and scalability in large-scale visual recognition tasks. Their ability to model non-local interactions makes them naturally attractive for multi-temporal reasoning. However, these advantages come with a significant computational burden. Standard self-attention requires quadratic complexity $\mathcal{O}(N^2)$ in both time and memory, considering $N$ as the number of tokens in the input. In multi-temporal BCD, each sample contains two high-resolution images, effectively doubling the sequence length and further amplifying the quadratic cost. As a result, existing Transformer-based approaches often rely on small patch sizes, heavy downsampling, or prohibitively small batch sizes, limiting their scalability and slowing experimentation in practical settings.

To alleviate this bottleneck, a wide range of \emph{efficient attention} mechanisms—kernelized linear attention \cite{Katharopoulos2020Linear,Choromanski2021Performer}, low-rank approximations \cite{Wang2020Linformer,Xiong2021Nystromformer}, and sparse or windowed attention widely used in vision \cite{Liu2021Swin,Wei2023Sparsifiner}—have been proposed to reduce complexity to near-linear in $N$. Despite their theoretical appeal, these methods often suffer from a critical issue: \emph{linear attention inherently collapses the rank of the key–value (KV) representation}, drastically reducing feature diversity and degrading accuracy in dense prediction tasks, especially in vision. Recent studies formalize this phenomenon and introduce the RALA \cite{Fan2025RALA} formulation to mitigate it by explicitly increasing the rank of the intermediate representations through learned token-wise weighting and output modulation. RALA breaks the traditional low-rank/accuracy trade-off observed in earlier linear attention variants.

Although RALA has demonstrated strong performance in natural image classification and general vision benchmarks \cite{Fan2025RALA}, its applicability to multi-temporal semantic segmentation remains unexplored. In particular, it is unknown whether the rank-preserving properties of RALA translate effectively to bi-temporal encoder–decoder architectures such as ChangeFormer, where accurate BCD requires both fine-grained local discrimination and global spatio–temporal reasoning. This gap is especially relevant given the computational constraints of VHR remote sensing, where efficient attention mechanisms must reduce complexity without degrading segmentation accuracy.

In this work, the gap is addressed by introducing \textbf{RALA-ChangeFormer}, an efficient variant of ChangeFormer in which each Transformer block replaces quadratic self-attention with the RALA mechanism, a rank-augmented linear attention formulation that expands the key–value space through a weighted mixture of multiple low-rank projections, preventing the collapse typically observed in standard linear attention. The goal is to reduce the computational and memory cost of multi-temporal BCD while preserving segmentation accuracy. The analysis examines how rank augmentation behaves within the hierarchical encoder–decoder design of ChangeFormer and quantifies its impact on performance, FLOPs, GPU memory consumption, and maximum feasible batch size.

The main contributions of this work are:
\begin{itemize}
    \item The introduction of \textbf{RALA-ChangeFormer}, a rank-augmented linear attention variant of ChangeFormer that achieves near-linear complexity while maintaining competitive accuracy in urban BCD.
    \item A detailed \textbf{computational analysis} quantifying reductions in GPU memory consumption, FLOPs, and the increase in feasible batch size enabled by rank-augmented attention.
    \item An extensive \textbf{experimental evaluation} on LEVIR-CD and DSIFN-CD demonstrating that RALA-ChangeFormer preserves segmentation accuracy while improving computational efficiency and training scalability.
\end{itemize}

This paper is organized as follows. Section~\ref{sec:related-work} reviews prior work on building change detection and efficient attention mechanisms. Section~\ref{sec:theoretical} introduces the theoretical foundations of ChangeFormer and RALA, including the key equations and mechanisms required for our integration. Section~\ref{sec:method} presents the proposed RALA-ChangeFormer method and describes how rank-augmented linear attention is incorporated into the architecture. Section~\ref{sec:experiments} details the experimental setup and analyzes the accuracy–efficiency trade-offs. Section~\ref{sec:conclusion} concludes the paper and outlines directions for future work.
